\section{Алгоритм CYK}
	\tikzsetfigurename{ClassicalParsing_CYK_}
\label{sec:lin_CYK}

Алгоритм CYK (Cocke-Younger-Kasami)~--- один из классических алгоритмов синтаксического анализа.
Его асимптотическая сложность в худшем случае~--- $O(n^3 \cdot |N|)$, где $n$~--- длина входной строки, а $N$~--- количество нетерминалов во входной грамматике~\sidecite{Hopcroft+Ullman/79/Introduction}.
\mytodo{Добавить псевдокод алгоритма CYK в окружении algorithm.}
\mytodo{Добавить схему доказательства корректности.}
\mytodo{При унификации примеров заменить грамматику и строку на сквозную (см.~\texttt{main.tex}).}

Для его применения необходимо, чтобы подаваемая на вход грамматика находилась в Нормальной Форме Хомского (НФХ)~\ref{sec:CNF}.
Других ограничений нет и, следовательно, данный алгоритм применим для работы с произвольными контекстно-свободными языками.
В основе алгоритма лежит принцип динамического программирования.
Используются два соображения:
\begin{enumerate}
    \item Из нетерминала $A$ выводится цепочка $\omega$ при помощи правила $A \to a$ тогда и только тогда, когда $a = \omega$:
          \[ A \derives \omega \iff \omega = a \]
    \item Из нетерминала $A$ выводится цепочка $\omega$ при помощи правила $A \to B C$ тогда и только тогда, когда существуют две цепочки $\omega_1$ и $\omega_2$ такие, что $\omega_1$ выводима из $B$, $\omega_2$ выводима из $C$ и при этом $\omega = \omega_1 \omega_2$:
          \[ A \derives[] B C \derives \omega \iff \exists \omega_1, \omega_2 : \omega = \omega_1 \omega_2, B \derives \omega_1, C \derives \omega_2 \]
          Переформулируем эти утверждения в терминах позиций в строке:
          \mytodo{Позиции в строке с 0 до $|\omega|$. Либо тут про позиции символа. Но тогда всё равно с 0 и до $|\omega|-1$}
          \[ A \derives[] B C \derives \omega \iff \exists k \in [1 \rng |\omega|] : B \derives \omega[1 \rng k], C \derives \omega[k+1 \rng |\omega|] \]
\end{enumerate}

В процессе работы алгоритма заполняется булева трехмерная матрица $M$ размера $n \times n \times  |N|$ таким образом, что
\[ M[i, j, A] = true \iff A \derives \omega[i \rng j]. \]

Первым шагом инициализируем матрицу, заполнив значения $M[i, i, A]$:
\begin{itemize}
    \item $M[i, i, A] = true$, если в грамматике есть правило $A \to \omega[i]$.\marginnote{TODO: Может надо заменить на 0 и 1?}
    \item $M[i, i, A] = false$, иначе.
\end{itemize}

Далее используем динамику: на шаге $m > 1$ предполагаем, что ячейки матрицы $M[i', j', A]$ заполнены для всех нетерминалов $A$ и пар $i', j': j' - i' < m$.
Тогда можно заполнить ячейки матрицы $M[i, j, A]$, где $j - i = m$ следующим образом:
\[ M[i, j, A] = \bigvee_{A \to B C} \bigvee_{k = i}^{j-1} M[i, k, B] \wedge M[k, j, C] \]

По итогу работы алгоритма значение в ячейке $M[0, |\omega|, S]$, где $S$~--- стартовый нетерминал грамматики, отвечает на вопрос о выводимости цепочки $\omega$ в грамматике.

\begin{example}\label{exmpl:CYK}
    \marginnote{TODO: Я бы сделал это подразделом, потому что пример всего один}
    Рассмотрим пример работы алгоритма CYK на грамматике правильных скобочных последовательностей в Нормальной Форме Хомского.
    \begin{align*}
        S   & \to A S_2 \mid \varepsilon & S_2 & \to b \mid B S_1 \mid S_1 S_3 & A & \to a \\
        S_1 & \to A S_2                  & S_3 & \to b \mid B S_1              & B & \to b
    \end{align*}

    Проверим выводимость цепочки $\omega = a a b b a b$.

    Так как трехмерные матрицы рисовать на двумерной бумаге не очень удобно, мы будем иллюстрировать работу алгоритма двумерными матрицами размера $n \times n$, где в ячейках указано множество нетерминалов, выводящих соответствующую подстроку.

    Шаг 1.
    Инициализируем матрицу элементами на главной диагонали:
    \marginnote{TODO: Я бы здесь написал снаружи нетерминалы, хотя бы в одной матричке}
    \[
        \begin{pmatrix}
            \{A\}       & \varnothing & \varnothing     & \varnothing     & \varnothing & \varnothing     \\
            \varnothing & \{A\}       & \varnothing     & \varnothing     & \varnothing & \varnothing     \\
            \varnothing & \varnothing & \{B, S_2, S_3\} & \varnothing     & \varnothing & \varnothing     \\
            \varnothing & \varnothing & \varnothing     & \{B, S_2, S_3\} & \varnothing & \varnothing     \\
            \varnothing & \varnothing & \varnothing     & \varnothing     & \{A\}       & \varnothing     \\
            \varnothing & \varnothing & \varnothing     & \varnothing     & \varnothing & \{B, S_2, S_3\}
        \end{pmatrix}
    \]

    Шаг 2.
    Заполняем диагональ, находящуюся над главной:
    \[
        \begin{pmatrix}
            \{A\}       & \varnothing & \varnothing                  & \varnothing     & \varnothing & \varnothing                  \\
            \varnothing & \{A\}       & \cellcolor{lightgray}\{S_1\} & \varnothing     & \varnothing & \varnothing                  \\
            \varnothing & \varnothing & \{B, S_2, S_3\}              & \varnothing     & \varnothing & \varnothing                  \\
            \varnothing & \varnothing & \varnothing                  & \{B, S_2, S_3\} & \varnothing & \varnothing                  \\
            \varnothing & \varnothing & \varnothing                  & \varnothing     & \{A\}       & \cellcolor{lightgray}\{S_1\} \\
            \varnothing & \varnothing & \varnothing                  & \varnothing     & \varnothing & \{B, S_2, S_3\}
        \end{pmatrix}
    \]
    В двух ячейках появились нетерминалы $S_1$ благодаря присутствию в грамматике правила $S_1 \to A S_2$.

    Шаг 3.
    Заполняем следующую диагональ:
    \[
        \begin{pmatrix}
            \{A\}       & \varnothing & \varnothing     & \varnothing                  & \varnothing & \varnothing                       \\
            \varnothing & \{A\}       & \{S_1\}         & \cellcolor{lightgray}\{S_2\} & \varnothing & \varnothing                       \\
            \varnothing & \varnothing & \{B, S_2, S_3\} & \varnothing                  & \varnothing & \varnothing                       \\
            \varnothing & \varnothing & \varnothing     & \{B, S_2, S_3\}              & \varnothing & \cellcolor{lightgray}\{S_2, S_3\} \\
            \varnothing & \varnothing & \varnothing     & \varnothing                  & \{A\}       & \{S_1\}                           \\
            \varnothing & \varnothing & \varnothing     & \varnothing                  & \varnothing & \{B, S_2, S_3\}
        \end{pmatrix}
    \]

    Шаг 4.
    И следующую за ней:
    \[
        \begin{pmatrix}
            \{A\}       & \varnothing & \varnothing     & \cellcolor{lightgray}\{S_1, S\} & \varnothing & \varnothing     \\
            \varnothing & \{A\}       & \{S_1\}         & \{S_2\}                         & \varnothing & \varnothing     \\
            \varnothing & \varnothing & \{B, S_2, S_3\} & \varnothing                     & \varnothing & \varnothing     \\
            \varnothing & \varnothing & \varnothing     & \{B, S_2, S_3\}                 & \varnothing & \{S_2, S_3\}    \\
            \varnothing & \varnothing & \varnothing     & \varnothing                     & \{A\}       & \{S_1\}         \\
            \varnothing & \varnothing & \varnothing     & \varnothing                     & \varnothing & \{B, S_2, S_3\}
        \end{pmatrix}
    \]

    Шаг 5.
    Заполняем предпоследнюю диагональ:
    \[
        \begin{pmatrix}
            \{A\}       & \varnothing & \varnothing     & \{S_1, S\}      & \varnothing & \varnothing                  \\
            \varnothing & \{A\}       & \{S_1\}         & \{S_2\}         & \varnothing & \cellcolor{lightgray}\{S_2\} \\
            \varnothing & \varnothing & \{B, S_2, S_3\} & \varnothing     & \varnothing & \varnothing                  \\
            \varnothing & \varnothing & \varnothing     & \{B, S_2, S_3\} & \varnothing & \{S_2, S_3\}                 \\
            \varnothing & \varnothing & \varnothing     & \varnothing     & \{A\}       & \{S_1\}                      \\
            \varnothing & \varnothing & \varnothing     & \varnothing     & \varnothing & \{B, S_2, S_3\}
        \end{pmatrix}
    \]

    Шаг 6.
    Завершаем выполнение алгоритма:
    \marginnote{TODO: Надо подумать как сделать так, чтобы не вылезало}
    \[
        \begin{pmatrix}
            \{A\}       & \varnothing & \varnothing     & \{S_1, S\}      & \varnothing & \cellcolor{lightgray}\{S_1, S\} \\
            \varnothing & \{A\}       & \{S_1\}         & \{S_2\}         & \varnothing & \{S_2\}                         \\
            \varnothing & \varnothing & \{B, S_2, S_3\} & \varnothing     & \varnothing & \varnothing                     \\
            \varnothing & \varnothing & \varnothing     & \{B, S_2, S_3\} & \varnothing & \{S_2, S_3\}                    \\
            \varnothing & \varnothing & \varnothing     & \varnothing     & \{A\}       & \{S_1\}                         \\
            \varnothing & \varnothing & \varnothing     & \varnothing     & \varnothing & \{B, S_2, S_3\}
        \end{pmatrix}
    \]

    Стартовый нетерминал находится в верхней правой ячейке, а значит цепочка $a a b b a b$ выводима в нашей грамматике.
\end{example}

\begin{example}
    Теперь выполним алгоритм на цепочке $\omega=abaa$.

    Шаг 1.
    Инициализируем таблицу:
    \[
        \begin{pmatrix}
            \{A\}       & \varnothing     & \varnothing & \varnothing \\
            \varnothing & \{B, S_2, S_3\} & \varnothing & \varnothing \\
            \varnothing & \varnothing     & \{A\}       & \varnothing \\
            \varnothing & \varnothing     & \varnothing & \{A\}
        \end{pmatrix}
    \]

    Шаг 2.
    Заполняем следующую диагональ:
    \[
        \begin{pmatrix}
            \{A\}       & \cellcolor{lightgray}\{S_1, S\} & \varnothing & \varnothing \\
            \varnothing & \{B, S_2, S_3\}                 & \varnothing & \varnothing \\
            \varnothing & \varnothing                     & \{A\}       & \varnothing \\
            \varnothing & \varnothing                     & \varnothing & \{A\}
        \end{pmatrix}
    \]

    Больше ни одну ячейку в таблице заполнить нельзя и при этом стартовый нетерминал отсутствует в правой верхней ячейке, а значит эта строка не выводится в грамматике правильных скобочных последовательностей.
\end{example}

Рассмотренный выше вариант алгоритма решает только задачу выводимости: в ячейках матрицы хранятся множества нетерминалов, а информация, необходимая для восстановления леса разбора, отбрасывается.
Модифицируем алгоритм так, чтобы по завершении работы можно было восстановить сжатое представление леса разбора (SPPF, раздел~\ref{sec:SPPF}) в варианте Рекерса (определение~\ref{def:basicSPPF}).

Основная идея состоит в том, чтобы хранить в ячейке матрицы не просто множество нетерминалов, а множество троек вида $(A, k, p)$, где $A$~--- нетерминал, $k$~--- точка расщепления, а $p$~--- номер продукции грамматики, приведённой к НФХ.
Клетка $M[i,j]$ матрицы, где $0 \leq i \leq j \leq n-1$, соответствует подстроке $w[i..j]$, то есть символам с индексами от $i$ до $j$ включительно.
Тройка $(A, k, p)$ в клетке $M[i,j]$ означает, что нетерминал $A$ выводит подстроку $w[i..j]$ при помощи продукции с номером $p$.
Для бинарной продукции $A \to B C$ точка расщепления $k$ ($i \leq k < j$) задаёт разбиение подстроки: $B$ выводит $w[i..k]$, а $C$ выводит $w[k+1..j]$.
Для терминальной продукции $A \to a$ точка расщепления не используется, и на её месте хранится индекс символа, то есть $k = i$.

Инициализация выполняется так же, как в исходном алгоритме: для каждой терминальной продукции $A \to a$ с номером $p$ и каждого индекса $i \in [0 \rng n-1]$ такого, что $w[i] = a$, в клетку $M[i,i]$ добавляется тройка $(A, i, p)$.

Динамика также повторяет исходный алгоритм, но фиксирует точку расщепления и номер продукции: для каждой бинарной продукции $A \to B C$ с номером $p$ и каждой точки расщепления $k$, $i \leq k < j$, если в клетке $M[i,k]$ найдётся тройка с нетерминалом $B$, а в клетке $M[k+1,j]$~--- тройка с нетерминалом $C$, то в клетку $M[i,j]$ добавляется тройка $(A, k, p)$.

По завершении основного прохода, если стартовый нетерминал $S$ присутствует в правой верхней клетке $M[0,n-1]$, то вывод успешен и SPPF восстанавливается рекурсивным проходом от корня.
В SPPF вершины нумеруются позициями между символами (определение~\ref{def:PositionInString}): клетке $M[i,j]$ соответствует символьная вершина $X_{i,j+1}$, а точке расщепления $k$~--- позиция $k+1$.
Для каждой тройки $(A, k, p)$ из клетки $M[i,j]$:
\begin{itemize}
    \item создаём (или переиспользуем уже созданную) символьную вершину $X_{i,j+1}$;
    \item создаём вершину-продукцию с номером $p$ и соединяем её с символьной вершиной $X_{i,j+1}$;
    \item если $p$~--- терминальная продукция $A \to a$, то создаём терминальный лист $a_{i,i+1}$ и соединяем с ним вершину-продукцию;
    \item если $p$~--- бинарная продукция $A \to B C$, то рекурсивно обрабатываем клетки $M[i,k]$ (для $B$) и $M[k+1,j]$ (для $C$), после чего соединяем вершину-продукцию с символьными вершинами $X_{i,k+1}$ и $X_{k+1,j+1}$.
\end{itemize}
Переиспользование символьных вершин по паре $(A, i, j)$ даёт разделение общих поддеревьев, а переиспользование вершин-продукций~--- упаковку альтернатив (см.~определение~\ref{def:basicSPPF}).
Заметим, что в матрице могут присутствовать тройки, не достигаемые при проходе от корня; в SPPF они не попадают.

\begin{example}\label{exmpl:CYK_SPPF}
    Проиллюстрируем работу модифицированного алгоритма на неоднозначной грамматике
    \begin{alignat*}{3}
    1) \ & S &&\rightarrow a \\
    2) \ & S &&\rightarrow S\ S \\
    3) \ & S &&\rightarrow S\ S\ S \\
    \end{alignat*}
    и цепочке $w = a a a$.
    Приведём грамматику к НФХ (раздел~\ref{sec:CNF}): продукцию $S \to S S S$ длины три заменим на две бинарные продукции, введя новый нетерминал $N_1$:
    \begin{alignat*}{3}
    1) \ & S &&\rightarrow a \\
    2) \ & S &&\rightarrow S\ S \\
    3) \ & S &&\rightarrow S\ N_1 \\
    4) \ & N_1 &&\rightarrow S\ S \\
    \end{alignat*}

    Длина цепочки равна $3$, поэтому заполняются клетки $M[i,j]$ для $0 \leq i \leq j \leq 2$; строка входа имеет вид
    \[
    \begin{pNiceMatrix}[margin=2pt] a & a & a \end{pNiceMatrix}
    \]

    Шаг 1.
    Для каждого индекса $i$ добавляем в клетку $M[i,i]$ тройку, отвечающую продукции $1)$ $S \to a$:
    \[
    \begin{pNiceMatrix}[color-inside,first-row,code-for-first-row = \arabic{jCol},first-col,code-for-first-col = \arabic{iRow}]
     &  &  &  \\
     & \cellcolor{yellow}{\{(S, 0, 1)\}} & \cdot & \cdot \\
     & \cdot & \cellcolor{yellow}{\{(S, 1, 1)\}} & \cdot \\
     & \cdot & \cdot & \cellcolor{yellow}{\{(S, 2, 1)\}} \\
    \end{pNiceMatrix}
    \]

    Шаг 2.
    Заполняем клетки, соответствующие подстрокам длины $2$.
    Для клетки $M[0,1]$ единственная точка расщепления~--- $k = 0$; так как и $M[0,0]$, и $M[1,1]$ содержат нетерминал $S$, по продукции $2)$ $S \to S S$ получаем тройку $(S, 0, 2)$, а по продукции $4)$ $N_1 \to S S$~--- тройку $(N_1, 0, 4)$.
    Аналогично для клетки $M[1,2]$ с точкой расщепления $k = 1$ получаем тройки $(S, 1, 2)$ и $(N_1, 1, 4)$:
    \[
    \begin{pNiceMatrix}[color-inside,first-row,code-for-first-row = \arabic{jCol},first-col,code-for-first-col = \arabic{iRow}]
     &  &  &  \\
     & \{(S, 0, 1)\} & \cellcolor{yellow}{\{(N_1, 0, 4), (S, 0, 2)\}} & \cdot \\
     & \cdot & \{(S, 1, 1)\} & \cellcolor{yellow}{\{(N_1, 1, 4), (S, 1, 2)\}} \\
     & \cdot & \cdot & \{(S, 2, 1)\} \\
    \end{pNiceMatrix}
    \]

    Шаг 3.
    Заполняем клетку $M[0,2]$, соответствующую всей цепочке.
    Для точки расщепления $k = 0$ используются клетки $M[0,0]$ и $M[1,2]$: по продукции $2)$ получаем тройку $(S, 0, 2)$, по продукции $3)$ $S \to S N_1$~--- тройку $(S, 0, 3)$, а по продукции $4)$~--- тройку $(N_1, 0, 4)$.
    Для точки расщепления $k = 1$ используются клетки $M[0,1]$ и $M[2,2]$: по продукции $2)$ получаем тройку $(S, 1, 2)$, а по продукции $4)$~--- тройку $(N_1, 1, 4)$:
    \begin{center}
    \resizebox{\textwidth}{!}{$\begin{pNiceMatrix}[color-inside,first-row,code-for-first-row = \arabic{jCol},first-col,code-for-first-col = \arabic{iRow}]
     &  &  &  \\
     & \{(S, 0, 1)\} & \{(N_1, 0, 4), (S, 0, 2)\} & \cellcolor{yellow}{\{(N_1, 0, 4), (N_1, 1, 4), (S, 0, 2), (S, 0, 3), (S, 1, 2)\}} \\
     & \cdot & \{(S, 1, 1)\} & \{(N_1, 1, 4), (S, 1, 2)\} \\
     & \cdot & \cdot & \{(S, 2, 1)\} \\
    \end{pNiceMatrix}$}
    \end{center}

    Стартовый нетерминал $S$ присутствует в правой верхней клетке $M[0,2]$, причём тремя различными способами, поэтому цепочка $a a a$ выводима.
    Восстановим SPPF проходом от корня (рис.~\ref{fig:sppf_cyk}).
    В нём символьные вершины помечены символом и диапазоном позиций (например, $S [0,3]$), а вершины-продукции~--- парой <<точка расщепления, номер продукции>>.
    У корня $S [0,3]$ три вершины-продукции~--- $(1,2)$, $(1,3)$ и $(2,2)$~--- по одной на каждый из трёх способов вывести всю цепочку.
    Обратим внимание на две особенности.
    Во-первых, символьные вершины $S [0,1]$, $S [1,2]$ и $S [2,3]$ переиспользуются: каждая из них имеет по три родительские вершины-продукции, что отражает неоднозначность грамматики~--- три дерева вывода цепочки $a a a$ <<упаковываются>> в один SPPF.
    Во-вторых, в клетке $M[0,2]$ присутствуют также тройки $(N_1, 0, 4)$ и $(N_1, 1, 4)$, однако соответствующая им символьная вершина $N_1 [0,3]$ недостижима из корня (стартовый нетерминал $S$ не обращается к $N_1$, покрывающему всю цепочку), поэтому в SPPF она не попадает.

    \begin{figure}
        \begin{center}
            \resizebox{0.8\textwidth}{!}{\begin{tikzpicture}
  \graph [layered layout, nodes={draw}, grow'=down, level sep=1.5cm, sibling sep=1.0cm] {
    n0 [as={S [0,3]}, fill=green!30];
    n1 [as={1, 2}];
    n2 [as={S [0,1]}];
    n3 [as={0, 1}];
    n4 [as={a\_\{0,1\}}];
    n5 [as={S [1,3]}];
    n6 [as={2, 2}];
    n7 [as={S [1,2]}];
    n8 [as={1, 1}];
    n9 [as={a\_\{1,2\}}];
    n10 [as={S [2,3]}];
    n11 [as={2, 1}];
    n12 [as={a\_\{2,3\}}];
    n13 [as={1, 3}];
    n14 [as={N\_1 [1,3]}];
    n15 [as={2, 4}];
    n16 [as={2, 2}];
    n17 [as={S [0,2]}];
    n18 [as={1, 2}];
    n0 -> n1;
    n0 -> n13;
    n0 -> n16;
    n1 -> n2;
    n1 -> n5;
    n2 -> n3;
    n3 -> n4;
    n5 -> n6;
    n6 -> n7;
    n6 -> n10;
    n7 -> n8;
    n8 -> n9;
    n10 -> n11;
    n11 -> n12;
    n13 -> n2;
    n13 -> n14;
    n14 -> n15;
    n15 -> n7;
    n15 -> n10;
    n16 -> n10;
    n16 -> n17;
    n17 -> n18;
    n18 -> n2;
    n18 -> n7;
  };
\end{tikzpicture}
}
        \end{center}
        \caption{SPPF для цепочки $a a a$ в неоднозначной грамматике $S \to a \mid S S \mid S S S$, построенный версией алгоритма CYK с сохранением троек.}
        \label{fig:sppf_cyk}
    \end{figure}
\end{example}
